\documentclass[12pt,letterpaper]{article}
\usepackage[utf8]{inputenc}
\usepackage[spanish]{babel}
\usepackage[margin=2.5cm]{geometry}
\usepackage{enumitem}
\usepackage{ulem}

\begin{document}

\begin{center}
    \textbf{UNIVERSIDAD DE ORIENTE}\\
    \textbf{NÚCLEO DE SUCRE}\\
    \textbf{ESCUELA DE CIENCIAS}\\
    \textbf{DEPARTAMENTO DE FÍSICA}\\
    \vspace{0.5cm}
    {\Large \textbf{\uline{FÍSICA PARA CIENCIAS DE LA SALUD}}}\\
\end{center}

\vspace{0.5cm}
\begin{flushright}
    CÓDIGO: 005-1713\\
    CRÉDITOS: 03\\
    \textbf{HORAS SEMANALES:}\\
    \textbf{2 TEORÍA}\\
    \textbf{4 PRÁCTICAS}
\end{flushright}

\vspace{0.5cm}

\noindent \textbf{SINTESIS DE CONOCIMIENTOS PREVIOS:}\\
Matemáticas generales e introducción al \textbf{cálculo diferencial e integral}.

\vspace{0.5cm}
\noindent \textbf{OBJETIVO GENERAL:}\\
Aplicar los conceptos básicos de la Física a la comprensión e interpretación de algunos fenómenos que tienen lugar en el cuerpo humano, así como a la comprensión del manejo y funcionamiento de instrumentos de utilidad en las ciencias de la salud.

\vspace{0.5cm}
\noindent \textbf{\uline{CONTENIDOS DE UNIDADES:}}

\vspace{0.3cm}
\noindent \uline{UNIDAD I:} Relación entre Física y las Ciencias de la Salud. Definición de ángulo plano. Sistemas de coordenadas. Vectores. Potencias de diez. Sistemas de unidades.

\vspace{0.3cm}
\noindent \uline{UNIDAD II:} Cinemática. Concepto de movimiento y sistemas de referencia. Movimiento rectilíneo uniforme. Movimiento uniformemente acelerado. Dinámica. Leyes de Newton. Efecto Fisiológico de las aceleraciones. Trabajo. Energía Cinética. Fuerzas conservativas. Energía Potencial. Potencia. 

\vspace{0.3cm}
\noindent \uline{UNIDAD III:} Fluidos. Densidad. Presión. Presión Hidrostática. Principios de Arquímedes y de Pascal. Presión atmosférica. Presión manométrica. Manómetros y la medida de la presión sanguínea. Fluidos en movimiento. Medidores de Venturi y de Pitot. Viscosidad. Flujo Sanguíneo y la Ecuación de Poiseuille. Resistencia Hemodinámica.

\vspace{0.3cm}
\noindent \uline{UNIDAD IV:} Líquidos. Calor de Vaporización. Tensión Superficial. Burbujas y tensión superficial en los pulmones. Capilaridad. Fenómenos de Transporte. Difusión. Primera y segunda ley de Fick. Osmosis y Osmosis inversa. Presión Osmótica.

\vspace{0.3cm}
\noindent \uline{UNIDAD V:} Electricidad. Fuerzas fundamentales. Ley de Coulomb. Campo Eléctrico. Energía Potencial electrostática. Potencial Eléctrico. Condensadores. Bioelectricidad. El potencial de Nernst. El potencial celular de equilibrio. La bomba de Potasio. Transmisión de impulsos nerviosos. La bomba de Sodio.

\vspace{0.5cm}
\noindent \textbf{\uline{BIBLIOGRAFÍA:}}
\begin{itemize}[label=\textbullet, leftmargin=*]
    \item Jerry D. Wilson, Anthony J. Buffa, Louis A., Bo Lou, \uline{Física}, 6\textsuperscript{da} Edición, Pearson Educación, México 2007.
\end{itemize}

\vspace{0.5cm}
\noindent Prof. César Rodríguez

\end{document}